\section{Introdução}
\label{sec:introducao}
O avanço dos modelos de linguagem (\textit{Large Language Models} -- LLMs) transformou o paradigma de interação entre humanos e máquinas. No entanto, a ausência inicial de padrões de comunicação entre LLMs e ferramentas externas gerou uma fragmentação técnica que motivou o surgimento do \textit{Model Context Protocol} (MCP). O protocolo vem sendo amplamente adotado como mecanismo de integração entre agentes inteligentes e ferramentas externas, alcançando aproximadamente 97 milhões de \textit{downloads} mensais de seus Kits de Desenvolvimento de Software (\textit{Software Development Kit} - SDK) no início de 2026 \cite{Avaya2026}. Seu principal diferencial reside no uso de ferramentas (\textit{tools}) descritas em linguagem natural para a execução de ações dinâmicas. Contudo, essa flexibilidade introduz uma dependência crítica da qualidade documental. Assim como \textit{code smells} comprometem a manutenibilidade de \textit{softwares} tradicionais, descrições ambíguas no contexto do MCP podem induzir agentes a erros e alucinações \cite{Taraghi2026}. Nesse contexto, alucinações correspondem à geração de comandos, referências ou conteúdos inexistentes, substituindo informações factuais por conteúdos produzidos a partir dos padrões aprendidos durante o treinamento do modelo \cite{Taraghi2026}.

Estudos recentes investigam a prevalência de anomalias (ou \textit{smells}) em descrições de ferramentas MCP por meio da \textbf{análise textual}. Essa abordagem avalia se a documentação comunica adequadamente o propósito, os parâmetros e as restrições de uso de uma ferramenta, verificando aspectos como clareza, completude e consistência sem considerar sua implementação técnica \cite{mcpAtFirstGlance}. Entretanto, evidências indicam que falhas relacionadas à configuração e à documentação representam 38,66\% dos problemas observados em servidores MCP \cite{Taraghi2026}. Considerando que a elaboração e manutenção dessas descrições dependem frequentemente de processos manuais desacoplados do código-fonte, existe o risco de surgirem discrepâncias semânticas entre a documentação e a implementação real. Diante desse cenário, o \textbf{problema central abordado neste estudo consiste na carência de investigações empíricas que confrontem sistematicamente as descrições de ferramentas MCP com sua implementação efetiva no código-fonte, visando identificar divergências que possam comprometer a confiabilidade de sistemas agênticos}.

A análise desse problema é relevante para a evolução sustentável de sistemas baseados em agentes de Inteligência Artificial (IA). Inconsistências estruturais podem reduzir a eficácia dos agentes e ampliar a superfície de exposição a vulnerabilidades, incluindo ataques de injeção de comandos \cite{Taraghi2026}. Além disso, investigações recentes identificaram a presença de anomalias em grande parte das ferramentas disponíveis no ecossistema MCP \cite{mcpAtFirstGlance}. Relatórios de segurança também alertam que ambiguidades documentais podem favorecer a execução de ações não autorizadas ou a exposição indevida de dados sensíveis \cite{NSA2026}. Nesse contexto, a incorporação de informações provenientes do código-fonte apresenta potencial para complementar abordagens tradicionais de avaliação de qualidade, permitindo confrontar a descrição textual com a implementação efetiva da ferramenta e fornecendo evidências adicionais sobre possíveis inconsistências.

Para investigar essa lacuna, este trabalho tem como \textbf{objetivo geral realizar uma avaliação empírica comparativa da qualidade de descrições de ferramentas MCP, integrando a análise textual tradicional ao contexto estrutural extraído do código-fonte}. Busca-se analisar em que medida uma abordagem de \textbf{avaliação enriquecida} pode contribuir para a identificação de discrepâncias semânticas e \textit{smells}, apoiando a construção de integrações mais confiáveis baseadas em modelos de linguagem. Para alcançar esse objetivo, definem-se os seguintes objetivos específicos: (i) mensurar o nível de clareza e completude das descrições em linguagem natural por meio de rubricas multidimensionais de qualidade fundamentadas na literatura \cite{2026arXiv260214878M}; (ii) quantificar o impacto da incorporação de propriedades estruturais extraídas do código-fonte na identificação de omissões e inconsistências técnicas não detectáveis exclusivamente por análise textual; e (iii) avaliar o potencial ganho de precisão diagnóstica proporcionado pelo contexto técnico, fornecendo subsídios estatísticos para a mitigação da dívida técnica e para o fortalecimento da robustez de sistemas baseados em MCP.

Como parte da investigação, será conduzida uma análise sobre um conjunto de servidores MCP de código aberto selecionados a partir de critérios de popularidade e maturidade do projeto. A avaliação se concentrará no confronto entre as descrições em linguagem natural das ferramentas e sua lógica de implementação, representada por informações estruturais extraídas do código-fonte. Espera-se que essa abordagem possibilite a identificação de inconsistências que não são facilmente observáveis por métodos exclusivamente textuais, como parâmetros implementados, mas não documentados, ou funcionalidades descritas sem correspondência evidente na implementação. Dessa forma, busca-se produzir evidências empíricas capazes de apoiar ou refutar a hipótese de que a incorporação do contexto técnico pode ampliar a capacidade de detecção de \textit{smells} em ferramentas MCP.

As próximas seções deste artigo estão organizadas da seguinte forma: a Seção 2 apresenta a Fundamentação Teórica sobre MCP, agentes baseados em LLMs e detecção de \textit{smells}. A Seção 3 discute os trabalhos relacionados. Por fim, a Seção 4 descreve os materiais e métodos empregados na condução do estudo.
